% Point de fin de séance à six du 09/11/2026 (action A8 du plan d'actions).
% Compiler depuis ce dossier : tectonic -o .. 2026-11-09_point-de-seance.tex
\documentclass[10pt,a4paper]{article}
\usepackage{style_cr}
\hypersetup{pdftitle={Point de séance du 09/11/2026}, pdfauthor={Reda El Mekkaoui},
            pdfsubject={Projet 20 Cerema, point de fin de séance}}
\newcommand{\groupe}[2]{{\small\bfseries\color{#1}\rule[-0.5pt]{2.4pt}{9pt}\enspace#2}\par\vspace{2pt}}

\begin{document}

\entete{PROJET 20 CEREMA \quad·\quad POINT DE FIN DE SÉANCE À SIX \quad·\quad 09/11/2026}
       {Point de séance du 09/11/2026}
       {Où en est chaque sous-groupe, et qui fait quoi pour la séance du jeudi 12/11.}
       {Participants : les six membres de l'équipe \quad·\quad Rédacteur : Reda El Mekkaoui \quad·\quad Prochain point : jeudi 12/11/2026, 8h30}

\begin{enbref}
\begin{itemize}
  \item \textbf{Plateforme} : le premier moteur (propulsion) est testé ; le deuxième (direction) se fait jeudi. Le module GPS et son antenne fonctionnent : ce qui bloque, c'est le branchement au Pixhawk, faute d'adaptateur 10 broches.
  \item \textbf{IA} : le modèle de détection de base analyse mal certaines images dégradées et hallucine sur les gouttes d'eau ; Guillaume va entraîner un nouveau modèle. La ré-identification a maintenant une mémoire des bateaux déjà vus.
  \item \textbf{Ensuite} : moteur de direction, puis configuration des missions avec le GPS, puis liaison Raspberry Pi–Pixhawk.
\end{itemize}
\end{enbref}

\section{Suivi des actions du 05/11}
\begin{tableau}{p{5.0cm} >{\bfseries}p{2.1cm} L}
\entetetab{\textbf{Action} & Qui & \textbf{État}}
GPS : dernière tentative sur le port GPS1 & Ewen & En cours, bloqué. Le module et l'antenne fonctionnent ; il manque un adaptateur 10 broches entre le module et le Pixhawk.\\\hline
Branchement des moteurs et des ESC & Rana & En cours : premier moteur testé, avec Selma.\\\hline
Tests de dégradation d'image & Guillaume & Fait. Résultat : le modèle de base analyse mal certaines images modifiées et hallucine sur les gouttes d'eau.\\\hline
Améliorer l'existant, en particulier les experts & Reda & En cours : recherche faite, tests jeudi.\\\hline
Ré-identification & Adam-Sabri & En cours : première version de la mémoire des bateaux déjà vus.\\\hline
Retester le modèle (« Base 4 ») & \mdseries\itshape à confirmer & Pas abordé sous ce nom. Les tests de Guillaume portent sur le modèle de détection de base.\\\hline
Télémétrie (Ewen), ArduPilot et QGroundControl (Selma), tests caméra (Guillaume), tests d'hallucination et petite boîte compacte (Reda) & \mdseries\itshape plusieurs & Non abordés ce jour. Selma a travaillé sur le moteur avec Rana. Point à faire jeudi à 8h30.\\\hline
\end{tableau}

\section{Qui fait quoi pour jeudi 12/11}
\groupe{cPlat}{Plateforme (lot L1)}
\begin{tableau}{>{\bfseries}p{2.1cm} p{4.0cm} L}
\entetetab{Qui & \textbf{Tâche} & \textbf{Détail}}
Rana & Moteur de direction & Brancher et tester le deuxième moteur, celui de la direction. Avec Selma.\\\hline
Selma & Moteur de direction & Avec Rana, comme pour le premier moteur.\\\hline
Ewen & GPS : brancher le module au Pixhawk & Trouver un adaptateur 10 broches. En attendant, emprunter le GPS de l'autre groupe si on en a besoin.\\\hline
Ewen & Configuration des missions & Dès qu'un GPS est branché, le nôtre ou celui de l'autre groupe.\\\hline
\end{tableau}

\needspace{8\baselineskip}
\vspace{0.6em}
\groupe{cIA}{IA (lot L2)}
\begin{tableau}{>{\bfseries}p{2.1cm} p{4.0cm} L}
\entetetab{Qui & \textbf{Tâche} & \textbf{Détail}}
Guillaume & Nouveau modèle de détection & L'entraîner davantage sur des images dégradées, gouttes d'eau en particulier, pour qu'il n'hallucine plus dessus.\\\hline
Reda & Tester les pistes d'amélioration des experts & Pistes trouvées depuis le 05/11. Les tester toutes.\\\hline
Adam-Sabri & Ré-identification & Suite non précisée en séance : à dire au point de 8h30.\\\hline
\end{tableau}

\section{Fait depuis la dernière séance}
\begin{itemize}
  \item \textbf{Plateforme} : Rana et Selma ont testé le premier moteur, celui de la propulsion. Ewen a vérifié que le module GPS et l'antenne fonctionnent ; seule la liaison entre le module et le Pixhawk ne marche pas.
  \item \textbf{IA} : Guillaume a testé le modèle de détection de base sur des images dégradées. Reda a cherché des pistes pour améliorer les experts. Adam-Sabri a ajouté à la ré-identification une mémoire des bateaux déjà vus :
  \begin{itemize}
    \item pour chaque bateau, on garde son immatriculation, son type et quelques signatures qui permettent de le reconnaître s'il sort de l'image puis réapparaît ;
    \item on sait ainsi qu'un bateau a déjà été vu, ce qui évite une approche inutile et permet de valider l'immatriculation plus vite ;
    \item on peut aussi repérer les incohérences.
  \end{itemize}
\end{itemize}

\section{Décisions}
\begin{enumerate}[label={\color{cExp}\bfseries\arabic*.}]
  \item \textbf{Ordre de travail de la plateforme} : moteur de direction, puis configuration des missions avec le GPS, puis liaison Raspberry Pi–Pixhawk, surtout une fois une première mission faite. Toute l'équipe est d'accord.
  \item \textbf{GPS en attendant} : tant que notre module n'est pas branché au Pixhawk, on emprunte le GPS de l'autre groupe quand on en a besoin.
\end{enumerate}

\section{Points ouverts et blocages}
\begin{itemize}
  \item \textbf{Adaptateur 10 broches} (module GPS vers Pixhawk) : on n'en a pas trouvé. Le Cerema n'en a pas non plus, et Ewen ne pense pas en trouver rapidement.
  \item \textbf{Si notre GPS ne marche pas au final} : on demandera probablement un autre module à Antoine Douchin. Rien n'est encore demandé.
  \item \textbf{Dépendance} : pour configurer les missions, il faut un GPS qui marche, le nôtre ou celui de l'autre groupe.
\end{itemize}

\section{Pour la suite}
\begin{itemize}
  \item \textbf{Plateforme} : après les missions, Ewen s'occupera de la liaison entre la Raspberry Pi et le Pixhawk. Elle a déjà été un peu regardée ; on s'y mettra surtout après une première mission.
\end{itemize}

\section{Suivi du plan d'actions}
\begin{tableau}{p{6.4cm} L}
\entetetab{\textbf{Indicateur} & \textbf{Ce jour}}
Tableau des tâches GitHub (A1) & Non abordé. En attendant, les tâches sont listées ici.\\\hline
Point de début de séance à 8h30 (A6) & Non relevé.\\\hline
Point de fin de séance (A7) & Fait : c'est ce compte rendu.\\\hline
Retards à l'arrivée ou après la pause (A10) & Non relevés.\\\hline
\end{tableau}

\redige{Rédigé par Reda El Mekkaoui d'après la transcription automatique, sans séparation fiable des voix : les interventions sont attribuées d'après l'ordre de parole (Guillaume, Reda, Adam-Sabri, Selma et Rana, Ewen). Corrections : « peak scope » → Pixhawk, « Serema » → Cerema, « ils ont fait la conclusion » → la propulsion (premier moteur). Une erreur ? Dites-le, on corrige.}

\end{document}
